% lambdastar03.tex -- round 479, PRIME.RDAGGER.LAMBDASTAR_03_KEYSTONE.01.
% REDUCED(kappa_high<=1.79e-2 at t*=8,N=5).  NOT lambda_*(0.3)>=0.
% scramble-sensitive (vacuous at L=0.3).  Companion: verify_lambdastar03.py.
% Discovery: experiments/tfpt-discovery/lambdastar03_probe.py.
% Fixed L=0.3 only.  NO RH CLAIM.  NO anti-RH claim.
\documentclass[11pt]{article}

\usepackage[a4paper,margin=2.45cm]{geometry}
\usepackage{amsmath,amssymb,amsthm}
\usepackage{booktabs,array,microtype,xcolor}
\usepackage[T1]{fontenc}
\usepackage[colorlinks=true,linkcolor=blue!50!black,urlcolor=blue!50!black]{hyperref}

\newtheorem{lemma}{Lemma}
\newtheorem{theorem}{Theorem}
\theoremstyle{remark}
\newtheorem{remark}{Remark}
\newcommand{\QW}{Q_{\mathrm W}}
\setlength{\emergencystretch}{2em}

\title{\bfseries Multiplier plus rank-one at $L=0.3$:\\[2mm]
a complete reduction, not yet $\lambda_*$\\[2mm]
\large PRIME.RDAGGER.LAMBDASTAR\_03\_KEYSTONE.01 --- round 479}
\author{Stefan Hamann}
\date{August 31, 2026}

\begin{document}
\maketitle

\begin{abstract}\noindent
At $L=0.3$ one has $P\equiv0$, so the compact-window Weil form
is a Fourier multiplier plus a rank-one update on each parity:
$\QW=A+\Pi$,
$A=(2\pi)^{-1}\int\sigma_A(t)|\hat h(t)|^2\,dt$,
$\Pi_{\mathrm{even}}=2\bigl|\langle h,\cosh(x/2)\rangle\bigr|^2$.
The zero $t_c$ of $\sigma_A$ and the elementary trace
$\mathrm{tr}(P_LB_{t_c}P_L)=2Lt_c/\pi$ are sealed.
The constant minorant $\sigma_A(0)K+\Pi$ is indefinite
(interlacing: rank one cannot lift a two-dimensional negative
space).  Actual $A$ on the even Fourier blocks $N=2..5$ is
positive definite (r478).  After five prolates of the band
$t_*=8$ the $K$-tail is positive with room
$c_{\mathrm{tail}}\ge0.2408$; low-band coupling is
$|\sigma_A(0)|\sqrt{\varepsilon}\le6.44\cdot10^{-3}$.
The split closes iff the high-band coupling of
$A_+=\sigma_A\,1_{|t|\ge8}$ satisfies
$\kappa_{\mathrm{high}}\le1.79\cdot10^{-2}$.
That last operator norm is the exact remainder.
Verdict \texttt{REDUCED($\kappa_{\mathrm{high}}\le1.79\cdot10^{-2}$ at $t_*=8,N=5$)}.
This $L$ lies in the Yoshida--Bombieri zone: the value is the
method, which must scale to $L=0.8$.  No RH claim.
\end{abstract}

\medskip\noindent
\textbf{Status.}\enspace
\texttt{REDUCED($\kappa_{\mathrm{high}}\le1.79\cdot10^{-2}$ at $t_*=8,N=5$)}.
No promotion of an RH claim.

\begin{center}
\fbox{\parbox{0.95\linewidth}{%
\textbf{Claim boundary.}  Nothing here is $\lambda_*(0.3)\ge0$.
The reduction is a theorem; the last coupling is open.
$L=0.3$ is inside the classical YB range $2L<\log2$:
this is a method-and-rigour upgrade, not a new zone.
No $L^*$ claim, no $R^\dagger$ claim, no RH claim,
no anti-RH claim.}}
\end{center}

\section{Convention}

$\mathrm{supp}\,h\subset[-L,L]$, $g=h*\widetilde h$,
$\QW=A-P+\Pi$, prime nodes $n\le e^{2L}$.
At $L=0.3$ the node set is empty, $P\equiv0$.
Plancherel $\hat h(t)=\int h(x)e^{itx}\,dx$,
$\|h\|_2^2=(2\pi)^{-1}\int|\hat h|^2$.
Symbol identity (r476, interval-matched):
$\sigma_A(t)=\mathrm{Re}\,\psi(\tfrac14+it/2)-\log\pi$,
even, increasing on $[0,\infty)$.
Closed value
$\sigma_A(0)=-(\gamma+\log\pi)-(\pi/2+3\log2)
\in[-5.372183419225667,\,-5.372183419225665]$.

\section{Multiplier plus rank one}

\begin{lemma}[Even--odd and $\Pi$]\label{lem:pi}
$A$ is a Fourier multiplier, hence
$A(\mathrm{even},\mathrm{odd})=0$.
For real $h$,
$\Pi(h)=\bigl|\langle h,e^{x/2}\rangle\bigr|^2
+\bigl|\langle h,e^{-x/2}\rangle\bigr|^2$,
so on even functions
$\Pi=2\bigl|\langle h,\cosh(x/2)\rangle\bigr|^2$
and on odd functions
$\Pi=2\bigl|\langle h,\sinh(x/2)\rangle\bigr|^2$.
Thus $\QW=Q_{\mathrm{even}}\oplus Q_{\mathrm{odd}}$.
Checked against $\mathtt{pole\_G}(0,2L)$: $\Pi(1)=16(\cosh L-1)$.
\end{lemma}

\begin{lemma}[The zero $t_c$]\label{lem:tc}
There is a unique $t_c>0$ with $\sigma_A(t_c)=0$, and
\[
 t_c\in\bigl[6.289835988836902,\,6.289835988836904\bigr]
\]
by the digamma enclosure (pad $10^{-16}$) plus the r476 identity.
Hence $\sigma_A(t)<0$ on $|t|<t_c$ and $\sigma_A(t)>0$ on $|t|>t_c$.
\end{lemma}

\section{The time-band operator and the trace ace}

Let $B_{t}$ be the Paley--Wiener projector
$(\widehat{Bh})(t)=1_{|t|<t}\hat h(t)$ and
$K_t=P_LB_tP_L$ on $L^2[-L,L]$.
The kernel is $\sin\bigl(t(x-y)\bigr)/\bigl(\pi(x-y)\bigr)$,
with on-diagonal value $t/\pi$.

\begin{lemma}[Trace]\label{lem:tr}
$\mathrm{tr}(K_t)=2Lt/\pi$ exactly.
At $t=t_c$,
$\mathrm{tr}(K_{t_c})\in[1.201270186632830,\,1.201270186632832]$.
\end{lemma}

Nystr\"om (Gauss--Legendre, $n=16$ vs $n=48$, top-five
difference $<10^{-14}$, trace match $10^{-8}$) gives the
Slepian numbers of $K_{t_c}$
\[
 \lambda=(0.859391,\,0.313778,\,0.027320,\,7.700\cdot10^{-4},\,
 1.120\cdot10^{-5},\,1.021\cdot10^{-7},\,\dots).
\]
Even indices $0,2,4$; odd $1,3,5$.  The first two even
eigenvalues already exhaust all but $1.12\cdot10^{-5}$ of the
even mass.  These $\lambda_n$ are Nystr\"om-stable working
constants, not a Lipschitz-interval enclosure.

\section{Why the constant minorant fails}

\begin{lemma}[Minorant, and why it is not enough]\label{lem:M}
$A\succeq\sigma_A(0)\,K_{t_c}$ as quadratic forms on $L^2[-L,L]$.
The rank-one lift $M=\sigma_A(0)K_{t_c}+2|\cosh\rangle\langle\cosh|$
has $\lambda_{\min}(M)\approx-3.448$ (Nystr\"om).
By Cauchy interlacing a rank-one PSD update lifts at most one
negative direction.  The second even prolate
($\lambda_2\approx0.0273$, overlap with $\cosh$ only $2.5\%$)
and the first odd prolate ($\lambda_1\approx0.314$,
$\|\sinh\|_2^2=\sinh L-L\approx0.00452$) remain negative
under $M$.  Therefore $M\succeq0$ is false, and cannot prove
$\QW\succeq0$.
\end{lemma}

The actual form $A$ is much larger than $\sigma_A(0)K$ on
every prolate with $\lambda_n\ll1$ (those functions live
outside the negative band).  One must use $A$ itself on the
essential block.

\section{What is already a SATZ}

\begin{theorem}[Even Fourier blocks, r478]\label{thm:block}
On $\mathrm{span}\{\cos(n\pi x/L):0\le n<N\}$ at $L=0.3$,
the Gram of $\QW$ is positive definite for $N=2,3,4,5$:
$N=2$ has $\lambda_{\min}\in[1.731281\cdot10^{-2},1.825107\cdot10^{-2}]$;
Higham $\mu_2\ge8.302\cdot10^{-3}$, $\mu_3\ge2.457\cdot10^{-3}$.
The first odd mode $\sin(\pi x/L)$ has $r\ge2.996973\cdot10^{-1}$.
\end{theorem}

This is $\lambda_*$ on a finite-dimensional subspace, not on
$L^2[-L,L]$.

\section{The reduction}

Write $\sigma_A(t)\ge\sigma_A(t_*)\,1_{|t|\ge t_*}+\sigma_A(0)\,1_{|t|<t_*}$
for any $t_*>t_c$ with $\sigma_A(t_*)>0$.  Then
\[
 A\ge\sigma_A(t_*)\,I+\bigl(\sigma_A(0)-\sigma_A(t_*)\bigr)K_{t_*}.
\]
Let $P_N$ be the projection onto the first $N$ prolates of
band $t_*$, $\varepsilon=\lambda_{N+1}(K_{t_*})$.
On $\mathrm{ran}(1-P_N)$,
\[
 A\ge\sigma_A(t_*)(1-\varepsilon)+\sigma_A(0)\varepsilon=:c_{\mathrm{tail}}.
\]
Because $K_{t_*}$ is diagonal in this basis,
$\langle K_{t_*}\,P_Nh,(1-P_N)h\rangle=0$, and the low-band
piece of $\langle A\,P_Nh,(1-P_N)h\rangle$ satisfies
\[
 \kappa_{\mathrm{low}}\le\bigl|\sigma_A(0)\bigr|\sqrt{\varepsilon}.
\]
The remaining high-band piece is
$\kappa_{\mathrm{high}}=\bigl\|P_NA_+ (1-P_N)\bigr\|$ with
$A_+=\sigma_A\,1_{|t|\ge t_*}$.

\begin{theorem}[Exact remainder]\label{thm:red}
Fix $t_*=8$ and $N=5$.  Sealed:
$\sigma_A(8)_{\mathrm{lo}}\ge0.240811$,
$\varepsilon=\lambda_6(K_8)\le1.4335\cdot10^{-6}$ (Nystr\"om-stable),
$c_{\mathrm{tail}}\ge0.240803$,
$\kappa_{\mathrm{low}}\le6.432\cdot10^{-3}$.
The even $3\times3$ Higham floor $\mu_3\ge2.457\cdot10^{-3}$
gives the numerical budget
\[
 \kappa_{\mathrm{high}}
 \;\le\;
 \sqrt{\mu_3\,c_{\mathrm{tail}}}-\kappa_{\mathrm{low}}
 \;\le\;1.79\cdot10^{-2}.
\]
If this operator-norm bound is proved, then
$\lambda_*(0.3)\ge c$ for an explicit $c>0$
(Schur of the even $3\times3$ against a positive tail;
the odd sector is strictly easier).
The bound is not proved in this round.
\end{theorem}

The r478 Hilbert--Schmidt majorant of the Fourier off-block
is $\sim0.27$ and overshoots the budget by an order of
magnitude.  That majorant is not forced: $A_+$ on the
$K$-tail is a different operator from the Fourier row-sum.

\section{Outlook at $L=0.8$}

There $P\not\equiv0$ with three nodes $n=2,3,4$.
As a quadratic form on $h$, $P$ is the Fourier multiplier
$\sigma_P(t)=\sum 2\Lambda(n)n^{-1/2}\cos(t\log n)$,
not a rank-three update of $A$.
$\Pi$ remains rank two (one per parity).
The pair $(A-P,\Pi)$ is still ``multiplier plus rank two'';
the negative set of $\sigma_A-\sigma_P$ is no longer a
single interval, so $t_c$ splits into a union of bands.
A four-dimensional Birman--Schwinger picture would be
accurate only after replacing $P$ by a finite-rank
surrogate.  The $L=0.3$ remainder
($\kappa_{\mathrm{high}}$ of a high-band multiplier)
is the piece that must stay explicit if the schema is to
scale.

\section{Sealed balance}

Probe \texttt{lambdastar03\_probe.py},
SPEC\_SHA \texttt{7a234db34c2ec6f0}.
Smoke $10/10$ in $0.36\,\mathrm{s}$; full $22/22$ in $0.44\,\mathrm{s}$.

\begin{center}
\small
\begin{tabular}{@{}llr@{}}
\toprule
object & status & constant\\
\midrule
$\sigma_A(0)$ & PROVED & $[-5.372183419225667,-5.372183419225665]$\\
$t_c$ & PROVED (digamma) & $[6.289835988836902,6.289835988836904]$\\
$\mathrm{tr}(K_{t_c})$ & PROVED & $[1.201270186632830,1.201270186632832]$\\
Slepian $\lambda_n(t_c)$ & Nystr\"om-stable & $0.859391,0.313778,0.027320,\dots$\\
$M=\sigma(0)K+\Pi$ & indefinite & $\lambda_{\min}\approx-3.448$\\
even $G^{(2)}$ $\lambda_{\min}$ & PROVED & $[1.731\cdot10^{-2},1.826\cdot10^{-2}]$\\
even $G^{(3)}$ Higham $\mu$ & PROVED & $\ge2.457\cdot10^{-3}$\\
odd $n=1$ Rayleigh & PROVED & $\ge2.997\cdot10^{-1}$\\
$c_{\mathrm{tail}}(8,5)$ & sealed & $\ge0.240803$\\
$\kappa_{\mathrm{low}}$ & PROVED & $\le6.432\cdot10^{-3}$\\
$\kappa_{\mathrm{high}}$ & OPEN & need $\le1.79\cdot10^{-2}$\\
\bottomrule
\end{tabular}
\end{center}

\section{Kill table}

\begin{center}
\small
\begin{tabular}{@{}lll@{}}
\toprule
gate & status & note\\
\midrule
$t_c$, trace, $\sigma_A(0)$ & PASS & every digit a pin\\
minorant $+$ rank one $\Rightarrow$ PD & KILL & interlacing, $M\prec0$\\
even/odd finite blocks & PASS & r478 $+$ odd $n=1$\\
tail positivity at $(8,5)$ & PASS & $c_{\mathrm{tail}}>0.24$\\
low-band coupling & PASS & $|\sigma(0)|\sqrt{\varepsilon}$\\
high-band $\kappa_{\mathrm{high}}$ & OPEN & the reduction\\
row-sum / HS Fourier & KILL & $\mathrm{HS}\sim0.27>0.018$\\
scramble / anti-list & PASS & $P=0$ vacuous; $e^{2L}$; no zeros\\
no RH / not $\lambda_*$ & PASS & type is REDUCED\\
\bottomrule
\end{tabular}
\end{center}

\section{What this does not do}

It does not prove $\lambda_*(0.3)\ge0$.  It does not enclose
the Slepian eigenvalues in a Lipschitz interval (Nystr\"om
stability only).  It does not treat $L>0.3$.  No cofinal
quantifier.  No RH claim.

\end{document}
